\documentclass{stylefiles/capstone}
\usepackage[utf8]{inputenc}
\usepackage[T1]{fontenc}
\usepackage{geometry}
\usepackage{graphicx}
\usepackage{float}
\usepackage{hyperref}
\usepackage[english,indonesian]{babel}

\title{Optimasi Pemrosesan Template LaTeX Menggunakan Pendekatan Berbasis AI}
\author{Academic Support Agent}
\date{\today}

\begin{document}
\maketitle

\begin{abstract}

\end{abstract}

\chapter{PENDAHULUAN}


\chapter{TINJAUAN PUSTAKA}


% Original content (commented out):
% \documentclass{stylefiles/capstone}

\begin{document}

\title{The Name of the Title is Hope}

%year
\capstoneyear{2020}
% options are Seminar, Project 1, Project 2
\capstonedocument{Seminar}
% options are Fall, Spring
\capstonesemester{Spring}

% For each author, provide the names and emails of the students as follows:
\author{Student Awesome}
\affiliation{%
  \institution{Computer Science, NYUAD}
}
\email{awesome@nyu.edu}

%\author{Student Eager}
%\affiliation{%
%  \institution{Computer Science, NYUAD}
%}
%\email{eager@nyu.edu}

%\author{Student Writer}
%\affiliation{%
%  \institution{Computer Science, NYUAD}
%}
%\email{writer@nyu.edu}

% Add your primary advisor's name, then comma-separate all secondary advisors' names.
% Please list advisors only and not secondary graders.
\advisors{Primary Advisor, Secondary Advisor}

\renewcommand{\shortauthors}{Awesome}

\begin{abstract}
  This is usually a short summary of the report. 1-2 paragraphs.
\end{abstract}


%Add a list of comma-separated keywords that are relevant to your report. They should be more specific than the list below
\keywords{database systems, human-computer interaction, NLP, cyber-security, computational geometry}



\maketitle

\section{Introduction}

\textbf{The section titles here and organization are merely suggestions. Please make sure that you use appropriate sections for your field of study.}

You need to learn latex commands to make sure your report is properly formatted. Points will be deducted for a poorly formatted document with incorrect references, figures or tables that do not fit within page margins, etc.

Pay attention to clarity and good writing style: The overall report should be well-written and persuasive with no spelling or grammar errors. Technical aspects of the proposal are clear and easy to understand. Writing style is simple: avoid dense sentences, flowery prose, passive voice, etc.


\begin{enumerate}
    \item Introduce the overall problem.
    \item Motivate the importance of the research. For example, why is this an
important problem? If you are building tools, what can these tools be used
for and by who?
    \item Provide the necessary background information on the problem.
    \item Explain the different challenges/problems. Why is this work hard/exciting?
\end{enumerate}

\section{Related Work}

You should cite at least 5 papers from top conferences/journals and describe
their core idea in relation to your proposal/report. Provide a summary of the techniques employed. You can also cite software tools if you did a quick test of them and report on their overall performance.

A sample reference is here ~\cite{chipapers}. Make sure you include proper bibtex by downloading them from dl.acm.org for example. You can refer to multiple papers as follows ~\cite{plpaper,sigraphpaper}

\section{Methodology}

\subsection{Methods}
If you implement any methods or design different algorithms, explain the
intuition behind them.
\begin{itemize}
    \item What is neat/intelligent about them?
    \item How do you expect them to perform?
    \item Why will they outperform other techniques?
    \item Feel free to include pseudo-code / algorithmic listings / hand-drawn sketches, etc.
    \item Expand on the ideas you proposed initially in the brainstorming posts. What
is neat/intelligent about them? How do you expect them to perform? Why
will they outperform other techniques (e.g., from a usability or performance
angle)?
\end{itemize}

If you provide proofs of correctness or formulaic derivations, list the proof
or the derivation and explain it in detail.

\section{Evaluation}
\begin{enumerate}
    \item How do you plan to evaluate your proposed methods? Is this a
thorough evaluation process? Would there be alternatives and why may they
not be applicable?
\item If you evaluate your preliminary methods/formulas, explain the evaluation
process/experimental setup
\item List the hypothesis and – possibly – preliminary results.
\item Provide appropriate visualizations of these results.
\end{enumerate}

\section{Project Timeline}

For the seminar, provide a two semester plan on what you intend to do and what you hope to achieve. For project 1, provide a timeline for project 2.
While this plan might change, it helps to think about a plan.

\section{Budget}

The budgetary requests are included in USD and are well justified. This section should only be included in the Seminar and Project 1.

\section{Conclusion}
Conclude your proposal emphasizing its main points and contributions.

\bibliographystyle{stylefiles/ACM-Reference-Format}
\bibliography{refs}

\end{document}
\endinput

\end{document}